\documentclass[10pt]{extarticle}
\usepackage[a4paper,margin=1.5cm]{geometry}
\usepackage[T1]{fontenc}
\usepackage{amsmath,amssymb}
\usepackage{graphicx}
\usepackage[export]{adjustbox}
\usepackage{multicol}
\usepackage{enumitem}
\usepackage{booktabs,array,tabularx}
\setlength{\columnsep}{0.7cm}
\setlength{\parindent}{0pt}
\setlength{\parskip}{3pt}
\pagestyle{empty}
\makeatletter
\newcommand{\qfigdef}[4]{\@namedef{qf@file@#1}{#2}\@namedef{qf@w@#1}{#3}\@namedef{qf@wide@#1}{#4}}
\newcommand{\qf@check}[1]{\@ifundefined{qf@file@#1}{\PackageError{score2}{Unknown figure id #1}{}}{}}
\newcommand{\qf@img}[4][]{\includegraphics[width=#2\dimexpr\@nameuse{qf@w@#4}pt\relax,max width=#3,max height=0.45\textheight,#1]{\@nameuse{qf@file@#4}}}
\newif\ifqf@full
\newif\ifqf@boxed
\newsavebox{\qf@box}
\newcommand{\qf@open}{\global\qf@fulltrue\par\global\qf@boxedtrue
  \setbox\qf@box\vbox\bgroup\hsize\textwidth\linewidth\textwidth}
\newcommand{\qf@emit}[1]{\par\egroup\global\qf@boxedfalse
  \ifdim\dimexpr\ht\qf@box+\dp\qf@box\relax<#1\textheight\noindent\usebox{\qf@box}\par\else\unvbox\qf@box\fi}
\newcommand{\qfullsplit}{\ifqf@boxed\qf@emit{0.85}\fi}
\newcommand{\qf@close}{\ifqf@boxed\qf@emit{0.45}\fi\global\qf@fullfalse}
\newcommand{\qfullbegin}{\end{multicols}\qf@open}
\newcommand{\qfullend}{\qf@close\begin{multicols}{2}}
\newcommand{\qfullstart}{\qf@open}
\newcommand{\qfulldone}{\qf@close}
\newcommand{\qfig}[2][1]{\qf@check{#2}%
  \ifqf@full\qf@blk{#1}{#2}\else
  \ifnum\@nameuse{qf@wide@#2}=1 \expandafter\@firstoftwo\else\expandafter\@secondoftwo\fi
  {\end{multicols}\begin{center}\qf@img{#1}{\textwidth}{#2}\end{center}\begin{multicols}{2}}%
  {\qf@blk{#1}{#2}}\fi}
\newcommand{\qf@blk}[2]{\par\smallskip{\centering\qf@img{#1}{\linewidth}{#2}\par}\smallskip}
\newcommand{\qfigi}[2][1]{\qf@check{#2}\qf@img[valign=c]{#1}{\linewidth}{#2}}
\makeatother
\newcommand{\papersec}[1]{\par\bigskip{\centering\Large\bfseries PAPER-#1\par}\medskip}
\newcommand{\partsec}[2]{\par\medskip\noindent{\large\bfseries PART-#1: #2}\par\smallskip}
\newcommand{\sect}[1]{\par\smallskip\noindent\textbf{#1}\par}
\newcommand{\qn}[1]{\par\smallskip\noindent\textbf{#1.}~}
\newcommand{\paraq}[2]{\par\smallskip\noindent\textbf{Paragraph for Questions #1 and #2}\par}
\newenvironment{qoptions}{\begin{itemize}[leftmargin=1.8em,itemsep=1pt,topsep=2pt,label={}]}{\end{itemize}}
\newcommand{\opt}[2]{\item[(#1)] #2}
% figure registry (generated by score2_tools.py finalize)
\qfigdef{p1_phy_s1_q01_1}{Score-II - Test 27 (23 Mar 2024)-img/p1_phy_s1_q01_1.png}{99.0}{0}
\qfigdef{p1_phy_s1_q09_1}{Score-II - Test 27 (23 Mar 2024)-img/p1_phy_s1_q09_1.png}{107.3}{0}
\qfigdef{p1_phy_s1_q10_1}{Score-II - Test 27 (23 Mar 2024)-img/p1_phy_s1_q10_1.png}{83.2}{0}
\qfigdef{p1_phy_s1_q16_1}{Score-II - Test 27 (23 Mar 2024)-img/p1_phy_s1_q16_1.png}{78.5}{0}
\qfigdef{p1_phy_s2_q02_1}{Score-II - Test 27 (23 Mar 2024)-img/p1_phy_s2_q02_1.png}{113.6}{0}
\qfigdef{p1_phy_s2_q03_1}{Score-II - Test 27 (23 Mar 2024)-img/p1_phy_s2_q03_1.png}{41.5}{0}
\qfigdef{p1_phy_s2_q04_1}{Score-II - Test 27 (23 Mar 2024)-img/p1_phy_s2_q04_1.png}{107.3}{0}
\qfigdef{p1_phy_s2_q05_1}{Score-II - Test 27 (23 Mar 2024)-img/p1_phy_s2_q05_1.png}{40.4}{0}
\qfigdef{p1_phy_s2_q06_1}{Score-II - Test 27 (23 Mar 2024)-img/p1_phy_s2_q06_1.png}{96.4}{0}
\qfigdef{p1_phy_s2_q08_1}{Score-II - Test 27 (23 Mar 2024)-img/p1_phy_s2_q08_1.png}{51.7}{0}
\qfigdef{p1_phy_s2_q09_1}{Score-II - Test 27 (23 Mar 2024)-img/p1_phy_s2_q09_1.png}{109.6}{0}
\qfigdef{p1_phy_s2_q10_1}{Score-II - Test 27 (23 Mar 2024)-img/p1_phy_s2_q10_1.png}{88.7}{0}
\qfigdef{p1_chem_s1_q02_1}{Score-II - Test 27 (23 Mar 2024)-img/p1_chem_s1_q02_1.png}{88.6}{0}
\qfigdef{p1_chem_s1_q06_1}{Score-II - Test 27 (23 Mar 2024)-img/p1_chem_s1_q06_1.png}{92.1}{0}
\qfigdef{p1_chem_s1_q12_1}{Score-II - Test 27 (23 Mar 2024)-img/p1_chem_s1_q12_1.png}{153.3}{0}
\qfigdef{p1_chem_s1_q13_1}{Score-II - Test 27 (23 Mar 2024)-img/p1_chem_s1_q13_1.png}{173.9}{0}
\qfigdef{p1_chem_s1_q13_2}{Score-II - Test 27 (23 Mar 2024)-img/p1_chem_s1_q13_2.png}{93.3}{0}
\qfigdef{p1_chem_s1_q13_3}{Score-II - Test 27 (23 Mar 2024)-img/p1_chem_s1_q13_3.png}{95.1}{0}
\qfigdef{p1_chem_s1_q13_4}{Score-II - Test 27 (23 Mar 2024)-img/p1_chem_s1_q13_4.png}{91.8}{0}
\qfigdef{p1_chem_s1_q13_5}{Score-II - Test 27 (23 Mar 2024)-img/p1_chem_s1_q13_5.png}{91.6}{0}
\qfigdef{p1_chem_s1_q15_1}{Score-II - Test 27 (23 Mar 2024)-img/p1_chem_s1_q15_1.png}{163.5}{0}
\qfigdef{p1_chem_s1_q15_2}{Score-II - Test 27 (23 Mar 2024)-img/p1_chem_s1_q15_2.png}{60.0}{0}
\qfigdef{p1_chem_s1_q15_6}{Score-II - Test 27 (23 Mar 2024)-img/p1_chem_s1_q15_6.png}{78.2}{0}
\qfigdef{p1_chem_s1_q15_3}{Score-II - Test 27 (23 Mar 2024)-img/p1_chem_s1_q15_3.png}{57.0}{0}
\qfigdef{p1_chem_s1_q15_4}{Score-II - Test 27 (23 Mar 2024)-img/p1_chem_s1_q15_4.png}{74.6}{0}
\qfigdef{p1_chem_s1_q16_1}{Score-II - Test 27 (23 Mar 2024)-img/p1_chem_s1_q16_1.png}{165.5}{0}
\qfigdef{p1_chem_s1_q16_2}{Score-II - Test 27 (23 Mar 2024)-img/p1_chem_s1_q16_2.png}{43.5}{0}
\qfigdef{p1_chem_s1_q16_3}{Score-II - Test 27 (23 Mar 2024)-img/p1_chem_s1_q16_3.png}{40.2}{0}
\qfigdef{p1_chem_s1_q16_4}{Score-II - Test 27 (23 Mar 2024)-img/p1_chem_s1_q16_4.png}{24.0}{0}
\qfigdef{p1_chem_s1_q16_5}{Score-II - Test 27 (23 Mar 2024)-img/p1_chem_s1_q16_5.png}{53.4}{0}
\qfigdef{p1_chem_s2_q07_1}{Score-II - Test 27 (23 Mar 2024)-img/p1_chem_s2_q07_1.png}{42.4}{0}
\qfigdef{p1_chem_s2_q08_1}{Score-II - Test 27 (23 Mar 2024)-img/p1_chem_s2_q08_1.png}{54.6}{0}
\qfigdef{p1_chem_s2_q08_2}{Score-II - Test 27 (23 Mar 2024)-img/p1_chem_s2_q08_2.png}{96.8}{0}
\qfigdef{p1_chem_s2_q09_2}{Score-II - Test 27 (23 Mar 2024)-img/p1_chem_s2_q09_2.png}{100.6}{0}
\makeatletter\@beginparpenalty=10000\makeatother
\begin{document}
\begin{center}{\Large\bfseries JEE (Main+Advanced) Score-II --- Test 27 (23 Mar 2024)}\end{center}
\begin{multicols}{2}
\raggedcolumns
\partsec{1}{PHYSICS}
\sect{Section-I}

\qn{1} The variation of the instantaneous current (I) and the instantaneous emf (E) in a circuit is as shown in figure. Which of the following statements is correct ?
\qfig{p1_phy_s1_q01_1}
\begin{qoptions}
\opt{A}{The voltage lags behind the current by $\pi/2$.}
\opt{B}{The voltage leads the current by $\pi/2$.}
\opt{C}{The voltage and the current are in phase.}
\opt{D}{The voltage leads the current by $\pi$.}
\end{qoptions}

\qn{2} \textbf{Assertion :} In $\alpha$-emission, the $\alpha$-particle carries the major share of the energy.\\
\textbf{and}\\
\textbf{Reason :} In $\alpha$-decay, the $\alpha$-particle being lighter has more kinetic energy.
\begin{qoptions}
\opt{A}{Assertion is True, Reason is True; Reason is a correct explanation for Assertion}
\opt{B}{Assertion is True, Reason is True; Reason is NOT a correct explanation for Assertion}
\opt{C}{Assertion is True, Reason is False}
\opt{D}{Assertion is False, Reason is True}
\end{qoptions}

\qn{3} A liquid of mass $2m$ and specific heat $c$ is heated to a temperature $2T$. Another liquid of mass $m$ and specific heat $2c$ is heated to a temperature $T$. If these two liquids are mixed, the resulting temperature of the mixture is
\begin{qoptions}
\opt{A}{$(2/3)\,T$}
\opt{B}{$(8/5)\,T$}
\opt{C}{$(3/5)\,T$}
\opt{D}{$(3/2)\,T$}
\end{qoptions}

\qn{4} A travelling microscope has a main scale division $\dfrac14$ mm. 60 vernier divisions coincide with 59 main scale divisions. What will be the reading if the main scale reads 8.45 cm and the 30$^{\text{th}}$ division of the vernier coincides with the main scale division ?
\begin{qoptions}
\opt{A}{8.5625 cm}
\opt{B}{8.4625 cm}
\opt{C}{8.588 cm}
\opt{D}{8.4675 cm}
\end{qoptions}

\qn{5} The filament current in the electron gun of a Coolidge tube is increased and also the potential difference used to accelerate the electron is increased. As a result in the emitted radiation :-
\begin{qoptions}
\opt{A}{The intensity decreases while the minimum wavelength increases.}
\opt{B}{The intensity increases while the minimum wavelength decreases.}
\opt{C}{The intensity as well as the minimum wavelength increases.}
\opt{D}{The intensity as well as the minimum wavelength decreases.}
\end{qoptions}

\qn{6} Two spherical nuclei have mass numbers 512 and 64 with their radii $R_1$ and $R_2$, respectively. The ratio $\dfrac{R_1}{R_2}$ is
\begin{qoptions}
\opt{A}{1.0}
\opt{B}{1.5}
\opt{C}{2.0}
\opt{D}{2.5}
\end{qoptions}

\qn{7} The mobility of electrons and holes in a sample of insentric germanium at room temperature are 0.88 and 0.12 m$^2$/V-s. If electron and holes densities each are equal to $5.0\times10^{19}/\mathrm{m^3}$ then conductivity of germanium will be (in ohm meter) :-
\begin{qoptions}
\opt{A}{2}
\opt{B}{4}
\opt{C}{6}
\opt{D}{8}
\end{qoptions}

\qn{8} An ideal gas has instantaneous pressure, Volume $v$ and temperature $T$ follows process\\
$P^2V^3=$ constant. Its volume expansion coefficient will be :-
\begin{qoptions}
\opt{A}{$\dfrac2T$}
\opt{B}{$-\dfrac2T$}
\opt{C}{$\dfrac3T$}
\opt{D}{$-\dfrac3T$}
\end{qoptions}

\qn{9} Six real batteries are connected in a circuit as shown in figure. Match the following :
\qfig{p1_phy_s1_q09_1}
\par\nobreak
\begin{center}
\begin{tabular}{|c|p{3.4cm}|c|c|}\hline
 & \multicolumn{1}{c|}{\textbf{Column-I}} & & \textbf{Column-II}\\\hline
(a) & Potential of point A & (p) & zero\\\hline
(b) & Potential of point B & (q) & 2 V\\\hline
(c) & Potential of point C & (r) & 4 V\\\hline
(d) & Potential of point D & (s) & 6 V\\\hline
\end{tabular}
\end{center}
\begin{qoptions}
\opt{A}{(a - p), (b - p), (c - p), (d - q)}
\opt{B}{(a - p), (b - p), (c - p), (d - s)}
\opt{C}{(a-q), (b - r), (c- s), (d - p)}
\opt{D}{None of these}
\end{qoptions}

\qn{10} Two insulating plates are both uniformly charged in such a way that the potential difference between them is $V_2-V_1=40$ V. (i.e., plate 2 is at a higher potential). The plates are separated by $d=0.1$ m and can be treated as infinitely large. An electron is released from rest on the inner surface of plate 1. What is its speed when it hits plate 2 ?\\
($e=1.6\times10^{-19}$ C, $m_0=9.11\times10^{-31}$ kg)
\qfig{p1_phy_s1_q10_1}
\begin{qoptions}
\opt{A}{$2.65\times10^6$ m/s}
\opt{B}{$3.75\times10^6$ m/s}
\opt{C}{$1.87\times10^6$ m/s}
\opt{D}{$32\times10^{-19}$ m/s}
\end{qoptions}

\qn{11} On investigation of light from three different stars A, B and C, it was found that in the spectrum of A the intensity of orange colour is maximum, in B the intensity of indigo colour is maximum and in C the intensity of green colour is maximum. From these observations it can be concluded that :-
\begin{qoptions}
\opt{A}{the temperature of A is maximum, B is minimum and C is intermediate.}
\opt{B}{the temperature of A is maximum, C is minimum and B is intermediate.}
\opt{C}{the temperature of B is maximum, A is minimum and C is intermediate.}
\opt{D}{the temperature of C is maximum, B is minimum and A is intermediate.}
\end{qoptions}

\qn{12} Choose the correct statements :-\\
(a) The density of electric lines of force at a point is independent of the magnitude of electric intensity vector $E$ at that point\\
(b) The density of electric lines of force at a point is proportional to the magnitude of electric intensity vector $E$ at that point\\
(c) Charge particle may or may not move along electric field lines when released at rest.\\
(d) Charge particle always moves along electric field lines when released at rest.
\begin{qoptions}
\opt{A}{b, c}
\opt{B}{c, d}
\opt{C}{b, d}
\opt{D}{a, d}
\end{qoptions}

\qn{13} A bulb is located on a wall. Its image is to be obtained on a parallel wall with the help of convex lens. If the distance between parallel walls is `$d$' then required focal length of lens placed in between the walls is :-
\begin{qoptions}
\opt{A}{Only $\dfrac d4$}
\opt{B}{Only $\dfrac d2$}
\opt{C}{More than $\dfrac d4$ but less than $\dfrac d2$}
\opt{D}{Less than or equal to $\dfrac d4$}
\end{qoptions}

\qn{14} When thick photosensitive surface is illuminated with light of wavelength $\lambda$, the stopping potential is $v$. When the same surface is illuminated by light of wavelength $4\lambda$, the stopping potential is $v/10$. Threshold wavelength for metallic surface is :-
\begin{qoptions}
\opt{A}{$\dfrac{4\lambda}{3}$}
\opt{B}{$4\lambda$}
\opt{C}{$\dfrac{8\lambda}{3}$}
\opt{D}{$6\lambda$}
\end{qoptions}

\qn{15} \textbf{Statement-1 :} Two particles moving in the same direction do not lose all their energy in a completely inelastic collision.\\
\textbf{Statement-2 :} Principle of conservation of momentum holds true for all kinds of collisions.
\begin{qoptions}
\opt{A}{Statement-1 is true, Statement-2 is false}
\opt{B}{Statement-1 is true, Statement-2 is true; Statement-2 is the correct explanation of Statement-1}
\opt{C}{Statement-1 is true, Statement-2 is true; Statement-2 is not the correct explanation of Statement-1}
\opt{D}{Statement-1 is false, Statement-2 is true}
\end{qoptions}

\qn{16} The diagram shows a force-extension graph for a rubber band. Consider the following statements-\\
I. It will be easier to compress this rubber than expand it.\\
II. Rubber does not return to its original length after it is stretched.\\
III. The rubber band will get heated if it is stretched and released.\\
Which of these can deduced from the graph-
\qfig{p1_phy_s1_q16_1}
\begin{qoptions}
\opt{A}{III only}
\opt{B}{II and III}
\opt{C}{I and III}
\opt{D}{I only}
\end{qoptions}

\qn{17} The intensity of light from a source is $\left(\dfrac{500}{\pi}\right)$ W/m$^2$. Find the amplitude of electric field in this wave.
\begin{qoptions}
\opt{A}{$\sqrt3\times10^2$ N/C}
\opt{B}{$2\sqrt3\times10^2$ N/C}
\opt{C}{$\dfrac{\sqrt3}{2}\times10^2$ N/C}
\opt{D}{$2\sqrt3\times10^1$ N/C}
\end{qoptions}

\qn{18} A barometer reads 75.0 cm on a steel scale. The room temperature is 30 $^\circ$C. The Scale is correctly graduated for 0$^\circ$C. The coefficient of linear expansion of steel is $\alpha=1.2\times10^{-5}/^\circ$C and the coefficient of volume expansion of mercury is $\gamma=1.8\times10^{-4}/^\circ$C. Correct atmospheric pressure (in cm of Hg) is closest to :-
\begin{qoptions}
\opt{A}{75}
\opt{B}{75.4}
\opt{C}{74}
\opt{D}{74.6}
\end{qoptions}

\qn{19} For a step-up transformer, the turns ratio is 3 and its efficiency is 0.75. The current flowing in the primary coil is 2A and the voltage applied to it is 100 V. Then the voltage and the current flowing in the secondary coil are ............. respectively :-
\begin{qoptions}
\opt{A}{150 V, 1.5 A}
\opt{B}{300 V, 0.5 A}
\opt{C}{300 V, 1.5 A}
\opt{D}{150 V, 0.5 A}
\end{qoptions}

\qn{20} In Young's double slit experiment, the distance $d$ between the slits S$_1$ and S$_2$ is 1.0 mm. What should be the width of each slit so as to obtain 5 maxima of the two slit interference pattern within the central maximum of the single slit diffraction pattern?
\begin{qoptions}
\opt{A}{0.2 mm}
\opt{B}{0.4 mm}
\opt{C}{12 cm}
\opt{D}{0.1 mm}
\end{qoptions}

\sect{Section-II}

\qn{1} In a circuit for finding the resistance of a galvanometer by half deflection method, a 20 V battery and a high resistance of 10 k$\Omega$ are used. The figure of merit of the galvanometer is 100 $\mu$A/division. In the absence of shunt resistance, the galvanometer produces a deflection of $\theta=15$ divisions when current flows in the circuit. The value of the shunt resistance that can cause the deflection of $\theta/2$, is $x$ ohm. Find value of $x$.

\qn{2} Figure shows infinitely long bent wires in the plane of paper and each carrying a current `$i$' with `O' as origin and $x$ and $y$ axis as shown, co-ordinates of points P \& Q are (2m, 0) and (4m, 0) respectively. Magnetic induction at P is given by $\dfrac{\alpha\mu_0 i}{32\pi}$. Fill the value of $\alpha$.
\qfig{p1_phy_s2_q02_1}

\qn{3} A block of mass 2 kg is suspended from a string, is released from a height of 20 cm as shown in figure then what will be the impulse when string just becomes tight (in N-s)? (Take $g=10\ \mathrm{m/s^2}$)
\qfig{p1_phy_s2_q03_1}

\qn{4} A 40 Henry inductor coil is connected to a 10 ohm resistance in series as shown in figure. The rate of dissipation of energy (joule's heat) across resistance is equal to the rate at which magnetic energy is stored in the inductor at time $t=y\ \ln2$ second, the value of $y$ is :
\qfig{p1_phy_s2_q04_1}

\qn{5} A mass $m=2$ kg hangs from the wheel of radius $R=2$m. When released from rest the mass falls through a height $h=4$m in 2 second. The moment of inertia (in kg m$^2$) of the wheel is ($g=10\ \mathrm{m/s^2}$) :-
\qfig{p1_phy_s2_q05_1}

\qn{6} Consider a cylinderical region of time varying magnetic field such that $\pi R^2\dfrac{dB}{dt}=24$ volt, where $\dfrac{dB}{dt}>0$. Points A and B are shown in space such that line AB subtends angle $60^\circ$ at centre. Calculate the value of induced emf along the straight line path AB (in volts)
\qfig{p1_phy_s2_q06_1}

\qn{7} A mixture of plane polarised and unpolarised light falls normally on a polarizing sheet. On rotating the polarizing sheet about the direction of the incident beam, the transmitted intensity varies by a factor of 4. The ratio of the intensities $I_P$ and $I_0$, respectively of the polarized and unpolarized components in the incident beam is $\dfrac\alpha2$. The value of $\alpha$ is :-

\qn{8} A small rod ABC is put in Liquid of Refractive index $\left(n=\dfrac98\right)$ making an angle $4^\mathrm{o}$ with vertical. If it is viewed paraxially from above, it will look like bent shaped ABC$'$. The angle of bending ($\angle$CBC$'$) is given by $\alpha^\mathrm{o}$. Find the value of $2\alpha$.
\qfig{p1_phy_s2_q08_1}

\qn{9} A particle of mas $m$ ($\ll M$) is projected with speed $V_o$ towards a sphere of mass $M$ and Radius $R$, as indicated in diagram. Particle just grazes the surface of sphere. If speed of projection of particle is $V_o=\sqrt{\dfrac{xGM}{yR}}$ then calculate value of $x+y$\\
($x$ and $y$ are coprime).
\qfig{p1_phy_s2_q09_1}

\qn{10} The figure shows a RC circuit with a parallel plate capacitor. Before switching on the circuit, plate A of the capacitor has a charge $-Q_0$ while plate B has no net charge. Now, at $t=0$, the circuit is switched on. How much time (in second) will elapse before the net charge on plate A becomes zero? (Given $C=1\,\mu$F, $Q_0=1$mC, $\varepsilon=1000$ V and $R=\dfrac{2\times10^6}{\ln3}\ \Omega$)
\qfig{p1_phy_s2_q10_1}

\partsec{2}{CHEMISTRY}
\sect{Section-I}

\qn{1} Select the correct option :
\begin{qoptions}
\opt{A}{$\Delta_{\text{combustion}}$H$^\circ$ of H(g) $=$ $\Delta_{\text{atomization}}$H$^\circ$ of H$_2$(g)}
\opt{B}{$\Delta_{\text{combustion}}$H$^\circ$ of graphite(s) $=$ $\Delta_{\text{atomization}}$H$^\circ$ of CO$_2$(g)}
\opt{C}{$\Delta_{\text{combustion}}$H$^\circ$ of C$_2$H$_4$(g) $=$ $\Delta_{\text{formation}}$H$^\circ$ of C$_2$H$_4$(g)}
\opt{D}{$\Delta_{\text{atomization}}$H$^\circ$ of H$_2$(g) $=$ Bond energy of H$_2$(g)}
\end{qoptions}

\qn{2} 2 moles of A$_2$ is present in container-I and 4 moles of B$_2$(g) is present in container-II at 300 K separated by a narrow tube of negligible volume closed by a valve as shown in figure. On opening the valve reaction given below attain equilibrium.\\
\centerline{A$_2$(g) + B$_2$(g) $\rightleftharpoons$ 2AB(g), K$_\mathrm{P}$ = 4}
\qfig{p1_chem_s1_q02_1}
Mole of B$_2$(g) in container-I at equilibrium would be
\begin{qoptions}
\opt{A}{1.33}
\opt{B}{1.5}
\opt{C}{2.0}
\opt{D}{2.67}
\end{qoptions}

\qn{3} The radius of the orbit of an electron in Hydrogen like atom is 4.5a$_0$, where a$_0$ is the Bohr radius. Its orbit angular momentum is $\dfrac{3h}{2\pi}$. It is given that $h$ is Plancks constant and $R$ is Rydberg constant. The possible longest wavelength, when the electron de$-$excites, is
\begin{qoptions}
\opt{A}{$\dfrac{9}{32R}$}
\opt{B}{$\dfrac{9}{16R}$}
\opt{C}{$\dfrac{9}{5R}$}
\opt{D}{$\dfrac{4}{3R}$}
\end{qoptions}

\qn{4} In which of the following reaction, O$_2$ is not formed as one of the product :
\begin{qoptions}
\opt{A}{MnO$_2$ + H$_2$SO$_4$ + H$_2$O$_2$ $\longrightarrow$}
\opt{B}{K$_3$[Fe(CN)$_6$] + KOH + H$_2$O$_2$ $\longrightarrow$}
\opt{C}{Cr$_2$(SO$_4$)$_3$ + H$_2$O$_2$ + NaOH $\longrightarrow$}
\opt{D}{H$_2$O$_2$ + MnO$_4^{-}$ + H$_2$SO$_4$ $\longrightarrow$}
\end{qoptions}

\qn{5} 1.6 gm CH$_4$(g) undergo combustion at constant volume at 300K and 2.5 kcal of heat is released. Heat of combustion of 1.6 gm CH$_4$(g) undergo combustion at constant pressure and same temperature will be:\\
(Given : R = 2 cal/mole-K)
\begin{qoptions}
\opt{A}{$-2380$ cal}
\opt{B}{$-2620$ cal}
\opt{C}{$-4500$ cal}
\opt{D}{2500 cal}
\end{qoptions}

\qn{6} 100 ml aqueous solution of HCl and CH$_3$COOH was conductometric titrated against 0.1 M NaOH aqueous solution.
\qfig{p1_chem_s1_q06_1}
pH of the solution when 125 ml of NaOH was added is :\\
(K$_\mathrm{a}$ of CH$_3$COOH = $10^{-5}$, log2 = 0.3, log3 = 0.48)
\begin{qoptions}
\opt{A}{4.7}
\opt{B}{5.0}
\opt{C}{5.48}
\opt{D}{8.52}
\end{qoptions}

\qn{7} \textbf{Statement-1 :} The number of reacting species (atoms, ions or molecules) taking part in an elementary reaction, which must collide simultaneously in order to bring about a chemical reaction is called the molecularity of a reaction.\\
\textbf{Statement-2 :} The molecularity of reaction of both complex and elementary reactions can be determined.
\begin{qoptions}
\opt{A}{Statement-1 is true, Statement-2 is true and Statement-2 is correct explanation for Statement-1}
\opt{B}{Statement-1 is true, Statement-2 is true and Statement-2 is NOT the correct explanation for Statement-1}
\opt{C}{Statement-1 is true, Statement-2 is false}
\opt{D}{Statement-1 is false, Statement-2 is true}
\end{qoptions}

\qn{8} Which of the following statement(s) is/are correct ?\\
(i) Electron Affinity : S$^{-}$(g) $>$ O$^{-}$(g)\\
(ii) Electron gain Enthalpy is highest for Neon in periodic table\\
(iii) Ionisation Energy is highest for Neon in periodic table\\
(iv) Second Ionisation Energy : Li (g) $>$ C(g)
\begin{qoptions}
\opt{A}{i, ii only}
\opt{B}{iii and iv only}
\opt{C}{i, ii and iv only}
\opt{D}{ii, iii and iv only}
\end{qoptions}

\qn{9} $\mathrm{Borax}\xrightarrow{\ \Delta\ }\mathrm{A}\xrightarrow{\ \Delta\ }\mathrm{B}+\mathrm{C}$\\
During borax bead test `C' combines with metal oxide to give coloured bead. Which of the following statement is correct for the above reaction and borax bead test ?\\
(i) One mol of A when dissolved in water produces an alkaline solution.\\
(ii) C combines with CoO to give red coloured bead of cobalt (II) metaborate.\\
(iii) Some atoms in Borax are involved in back bonding\\
(iv) B is sodium metaborate
\begin{qoptions}
\opt{A}{i, iii only}
\opt{B}{i, iii and iv only}
\opt{C}{i and ii only}
\opt{D}{ii, iii and iv only}
\end{qoptions}

\qn{10} In bis-(dimethylglyoximato)nickel(II), the sum of number of Ni-N, Ni-O and Intramolecular Hydrogen bond are :
\begin{qoptions}
\opt{A}{4}
\opt{B}{5}
\opt{C}{6}
\opt{D}{7}
\end{qoptions}

\qn{11} \textbf{Assertion (A) :} Cr$^{2+}$ is oxidizing and and Mn$^{3+}$ is reducing although both have d$^4$ configuration.\\
\textbf{Reason (R) :} Cr$^{2+}$ configuration changes to acquire half filled t$_{2g}$ where as Mn$^{3+}$ configuration changes to acquire half filled d$^5$ configuration.
\begin{qoptions}
\opt{A}{Both A and R are true and R is the correct explanation of A.}
\opt{B}{Both A and R are true and R is NOT the correct explanation of A}
\opt{C}{A is true but R is false}
\opt{D}{A is false but R is true.}
\end{qoptions}

\qn{12} \qfig{p1_chem_s1_q12_1}
Which of the following is correct decreasing order of acidic strength ?
\begin{qoptions}
\opt{A}{B $>$ A $>$ D $>$ C}
\opt{B}{D $>$ A $>$ B $>$ C}
\opt{C}{A $>$ B $>$ C $>$ D}
\opt{D}{D $>$ B $>$ C $>$ A}
\end{qoptions}

\qn{13} \qfig{p1_chem_s1_q13_1}
The major products B and C are respectively :
\begin{qoptions}
\opt{A}{\qfigi{p1_chem_s1_q13_2}}
\opt{B}{\qfigi{p1_chem_s1_q13_3}}
\opt{C}{\qfigi{p1_chem_s1_q13_4}}
\opt{D}{\qfigi{p1_chem_s1_q13_5}}
\end{qoptions}

\qn{14}
\par\nobreak
\begin{center}
\begin{tabular}{|c|p{3.3cm}|c|p{2.0cm}|}\hline
 & \multicolumn{1}{c|}{\textbf{Column-I}} & & \multicolumn{1}{c|}{\textbf{Column-II}}\\\hline
(P) & CAN TEST & (i) & Scarlet Red\\\hline
(Q) & Molisch's Test & (ii) & Rhumann's Purple\\\hline
(R) & Ninhydrin Test & (iii) & Red\\\hline
(S) & $\beta$-naphthol (Azo dye test) & (iv) & Violet\\\hline
\end{tabular}
\end{center}
\begin{qoptions}
\opt{A}{P - iv, Q - ii, R - iii, S - i}
\opt{B}{P - iii, Q - iv, R - i, S - ii}
\opt{C}{P - i, Q - ii, R - ii, S - iv}
\opt{D}{P - iii, Q - iv, R - ii, S - i}
\end{qoptions}

\qn{15} \qfig{p1_chem_s1_q15_1}
The correct structure of dye is :
\begin{qoptions}
\opt{A}{\qfigi{p1_chem_s1_q15_2}}
\opt{B}{\qfigi{p1_chem_s1_q15_6}}
\opt{C}{\qfigi{p1_chem_s1_q15_3}}
\opt{D}{\qfigi{p1_chem_s1_q15_4}}
\end{qoptions}

\qn{16} \qfig{p1_chem_s1_q16_1}
The major product B will be :
\begin{qoptions}
\opt{A}{\qfigi{p1_chem_s1_q16_2}}
\opt{B}{\qfigi{p1_chem_s1_q16_3}}
\opt{C}{\qfigi{p1_chem_s1_q16_4}}
\opt{D}{\qfigi{p1_chem_s1_q16_5}}
\end{qoptions}

\qn{17} Given below are two statements : one is labelled as Assertion A and the other is labelled as Reason R :\\
\textbf{Assertion A :} Aryl halides cannot be prepared by replacement of hydroxyl group of phenol by halogen atom.\\
\textbf{Reason R :} Phenols react with halogen acids violently. In the light of the above statements, choose the most appropriate from the options given below:
\begin{qoptions}
\opt{A}{Both A and R are true but R is NOT the correct explanation of A}
\opt{B}{A is false but R is true}
\opt{C}{A is true but R is false}
\opt{D}{Both A and R are true and R is the correct explanation of A}
\end{qoptions}

\qn{18} Which among the following purification methods is based on the principle of ``Solubility'' in two different solvents?
\begin{qoptions}
\opt{A}{Column Chromatography}
\opt{B}{Sublimation}
\opt{C}{Distillation}
\opt{D}{Differential Extraction}
\end{qoptions}

\qn{19} Which of the following is the method of preparation of oxygen?
\begin{qoptions}
\opt{A}{2KClO$_3$ $\xrightarrow[\mathrm{MnO_2}]{\text{heat}}$ 2KCl + 3O$_2$}
\opt{B}{2Ag$_2$O(s) $\to$ 4Ag(s) + O$_2$(g)}
\opt{C}{Catalytic decomposition of hydrogen peroxide}
\opt{D}{All of these}
\end{qoptions}

\qn{20} Select the incorrect statement.
\begin{qoptions}
\opt{A}{There is gradual change from basic V$_2$O$_3$ to less basic V$_2$O$_4$ and to amphoteric V$_2$O$_5$}
\opt{B}{V$_2$O$_4$ dissolves in acids to give VO$^{2+}$ salts}
\opt{C}{V$_2$O$_5$ reacts with alkalies as well as acids to give V O$_4^{+}$ and V O$_4^{3-}$ respectively}
\opt{D}{V$_2$O$_5$ is amphoteric though mainly acidic}
\end{qoptions}

\sect{Section-II}

\qn{1} Following informations are given for the formation of octahedral complexes of Co$^{2+}$ in aqueous medium at 27$^\circ$C :\\
Co$^{2+}$(aq) + 6NH$_3$(aq) $\rightleftharpoons$ [Co(NH$_3$)$_6^{2+}$](aq) ;\\
K$_{\mathrm{eq1}}=e^{20}$ ; $\Delta$H$^0_1=-x$ kcal/mole,\\
$\Delta$S$^0_1=-30$ cal/K-mole\\
Co$^{2+}$(aq) + 6MeNH$_2$(aq)\\
$\rightleftharpoons$ [Co(MeNH$_2$)$_6^{2+}$](aq) ; K$_{\mathrm{eq2}}=e^{30}$ ;\\
$\Delta$H$^0_2=-2x$ Kcal/mole, $\Delta$S$^0_2=-y$ cal/K-mole\\
The value of $y$ is \underline{\hspace{2em}}.

\qn{2} When 1M Pb(NO$_3$)$_2$ solution is tirated with 1 M KI solution then what will be the osmotic pressure (in atm) when equivalence point was reached at 300K. (R = 0.08 atm-L/mole-K)

\qn{3} Standard reduction potential of reaction :\\
PbCl$_2$ + 2e$^{-}$ $\longrightarrow$ Pb + 2Cl$^{-}$ Is ($-$ ...........) $\times10^{-2}$ volt.\\
(Given : $\mathrm{E^{o}_{Pb^{2+}|Pb}}=-0.13$V , K$_{\mathrm{SP}}$ of PbCl$_2$ = $1\times10^{-5}$)

\qn{4} SO$_4^{2-}$ , XeO$_6^{4-}$ , CO$_3^{2-}$ , HCOO$^{-}$ , NO$_3^{-}$ , CO$_2$\\
$x$ = number of molecules /Ions where resonance takes place\\
$y$ = No. of Ions where X - O bond order is greater than or equal to 1.5 find the value of $x+y$.

\qn{5} Diborane + Ammonia\\
\centerline{$\xrightarrow[\text{50\% Efficiency}]{\text{Room temperature}}$ A $\xrightarrow[\text{75\% Efficiency}]{\text{Heat}}$ B + C(gas)}
Find the mass of B (in gm) produced from 4 moles of Diborane. (B = 11, N = 14)

\qn{6} For how many of the following the colour is correctly matched with the ppt./residue etc?\\
A. [Fe(OH)$_2$(CH$_3$COO$^{-}$)] - Brown Red\\
B. PbCrO$_4$ $-$ Red\\
C. CrO$_5$ - Yellow\\
D. ZnO (Hot) $-$ White\\
E. Cu$_2$[Fe(CN)$_6$] - Brown\\
F. Fe$_4$[Fe(CN)$_6$]$_3$ - Prussian Blue\\
G. K$_3$[Co(NO$_2$)$_6$] - Red

\qn{7} Maximum number of moles of hot conc. HI that can be consumed by 1 mole of
\par\nobreak
\begin{center}\qfigi{p1_chem_s2_q07_1}\end{center}

\qn{8} \begin{center}\qfigi{p1_chem_s2_q08_1}\qfigi{p1_chem_s2_q08_2}\end{center}
then $(X+Y)$ is :

\qn{9} The number of Cannizaro reaction(s) that occurs in the given transformation is :\\
CH$_3$--CHO + 4 H--CHO
\qfig{p1_chem_s2_q09_2}

\qn{10} $\mathrm{Xe(g)+3F_2(g)}\ \xrightarrow[\text{60--70 bar}]{\text{573K}}\ \mathrm{XeF_6(s)}$ \quad (1 : $x$ ratio)\\
The value of $x$ is :-

\partsec{3}{MATHEMATICS}
\sect{Section-I}

\qn{1} In a survey, it was found out that 100 people don't use any of laptop, mobile or wrist watches. 80 persons use all the three gadgets. There are 150 who use laptop and mobile, 200 who use mobile and wristwatch and 200 who use laptop and wristwatch. The number of people who use only laptop, only mobile and only wristwatch is equal. If this survey was conducted on 1000 persons, How many people use only wristwatch ?
\begin{qoptions}
\opt{A}{200}
\opt{B}{170}
\opt{C}{150}
\opt{D}{900}
\end{qoptions}

\qn{2} If $R=\{(a,b):a,b\in R,\ 1+ab>0\}$ is a relation on set of real numbers $R$, then $R$ is
\begin{qoptions}
\opt{A}{Equivalence}
\opt{B}{Symmetric and transitive but not reflexive.}
\opt{C}{Reflexive and symmetric but not transitive}
\opt{D}{Neither symmetric nor transitive}
\end{qoptions}

\qn{3} \textbf{Given :} A circle, $2x^2+2y^2=5$ and a parabola, $y^2=4\sqrt5x$.\\
\textbf{Assertion A :} An equation of a common tangent to these curves is $y=x+\sqrt5$.\\
\textbf{Reason R :} If the line, $y=mx+\dfrac{\sqrt5}{m}$ ($m\neq0$) is their common tangent, then $m$ satisfies $m^4-3m^2-4=0$
\begin{qoptions}
\opt{A}{Both \textbf{A} and \textbf{R} are correct and \textbf{R} is the correct explanation of \textbf{A}}
\opt{B}{Both \textbf{A} and \textbf{R} are correct but \textbf{R} is \textbf{NOT} the correct explanation of \textbf{A}}
\opt{C}{\textbf{A} is correct but \textbf{R} is not correct}
\opt{D}{\textbf{A} is not correct but \textbf{R} is correct}
\end{qoptions}

\qn{4} If $z+\sqrt2|z+1|+i=0$, then $|z|^2$ is
\begin{qoptions}
\opt{A}{5}
\opt{B}{6}
\opt{C}{7}
\opt{D}{8}
\end{qoptions}

\qn{5} Let $\lambda\neq0$ be a real number. Let $\alpha,\beta$ be the roots of the equation $14x^2-31x+3\lambda=0$ and $\alpha,\gamma$ the roots of the equation $35x^2-53x+4\lambda=0$. Then $\dfrac{3\alpha}{\beta}$ and $\dfrac{4\alpha}{\gamma}$ are the roots of the equation
\begin{qoptions}
\opt{A}{$7x^2-245x+250=0$}
\opt{B}{$49x^2-245x+250=0$}
\opt{C}{$49x^2+245x+250=0$}
\opt{D}{$7x^2+245x-250=0$}
\end{qoptions}

\qn{6} If $\vec a$ and $\vec b$ are two perpendicular unit vectors such that $\vec x=\vec b-(\vec a\times\vec x)$, then $|\vec x|=$
\begin{qoptions}
\opt{A}{1}
\opt{B}{$\sqrt2$}
\opt{C}{$\dfrac1{\sqrt2}$}
\opt{D}{$\sqrt3$}
\end{qoptions}

\qn{7} The derivative of the function\\
$f(x)=\cos^{-1}\left\{\dfrac1{\sqrt{13}}(2\cos x-3\sin x)\right\}$\\$\phantom{f(x)=}+\sin^{-1}\left\{\dfrac1{\sqrt{13}}(2\sin x+3\cos x)\right\}$\\
with respect to $\sqrt{1+x^2}$ is $\left(x\in\left(0,\dfrac\pi6\right)\right)$
\begin{qoptions}
\opt{A}{$2x$}
\opt{B}{$2\sqrt{1+x^2}$}
\opt{C}{$\dfrac2x\sqrt{1+x^2}$}
\opt{D}{$\dfrac{2x}{\sqrt{1+x^2}}$}
\end{qoptions}

\qn{8} Let $f(x)=\begin{cases}\left|1-2x^2\right|, & 0\leqslant x<1\\ \left[x^2-2x\right], & 1\leqslant x<2\end{cases}$\\
If $m,n$ are number of points of discontinuity and non-differentiability of $f(x)$ in $(0,2)$ then ($[.]$ represents greatest integer function)
\begin{qoptions}
\opt{A}{$m=1,\ n=2$}
\opt{B}{$m=1,\ n=1$}
\opt{C}{$m=2,\ n=2$}
\opt{D}{None of these}
\end{qoptions}

\qn{9}
\par\nobreak
\begin{center}\small
\begin{tabular}{|c|p{3.6cm}|c|c|}\hline
 & \multicolumn{1}{c|}{\textbf{List-I}} & & \textbf{List-II}\\\hline
(A) & Suppose $F(n+1)=\dfrac{2F(n)+1}{2}$ for $n=1,2,3,\dots$ and $F(1)=2$. Then $F(101)$ equals to \dots\dots & (I) & 42\\\hline
(B) & If $a_1,a_2,a_3\dots a_{21}$, are in A.P. and $a_3+a_5+a_{11}+a_{17}+a_{19}=10$, then the value of $\displaystyle\sum_{i=1}^{21}a_i$ is \dots.. & (II) & 1620\\\hline
(C) & $10^{\text{th}}$ term of the sequence $1+5+13+29+\dots$ is \dots.. & (III) & 52\\\hline
(D) & The sum of all two digit numbers which are neither divisible by 2 nor by 3 is \dots\dots & (IV) & 2045\\\hline
\end{tabular}
\end{center}
\begin{qoptions}
\opt{A}{(A) $\to$ III, (B) $\to$ II, (C) $\to$ IV, (D) $\to$ I}
\opt{B}{(A) $\to$ III, (B) $\to$ I, (C) $\to$ IV, (D) $\to$ II}
\opt{C}{(A) $\to$ II, (B) $\to$ I, (C) $\to$ IV, (D) $\to$ I}
\opt{D}{(A) $\to$ III, (B) $\to$ I, (C) $\to$ II, (D) $\to$ IV}
\end{qoptions}

\qn{10} There are 10 coins out of which one is double headed. A coin is chosen at random and tossed 3 times. Head appears on all the 3 times. The probability that chosen coin is a double headed coin is
\begin{qoptions}
\opt{A}{$\dfrac12$}
\opt{B}{$\dfrac3{17}$}
\opt{C}{$\dfrac8{17}$}
\opt{D}{$\dfrac{11}{17}$}
\end{qoptions}

\qn{11} The area of the region between the curves $y=\sqrt{\dfrac{1+\sin x}{\cos x}}$ and $y=\sqrt{\dfrac{1-\sin x}{\cos x}}$ bounded by the lines $x=0$ and $x=\dfrac\pi4$ is
\begin{qoptions}
\opt{A}{$\displaystyle\int_{\sqrt2+1}^{\sqrt2-1}\dfrac{t}{(1+t^2)\sqrt{1-t^2}}\,dt$}
\opt{B}{$\displaystyle\int_0^{\sqrt2-1}\dfrac{4t}{(1+t^2)\sqrt{1-t^2}}\,dt$}
\opt{C}{$\displaystyle\int_0^{\sqrt2+1}\dfrac{4t}{(1+t^2)\sqrt{1-t^2}}\,dt$}
\opt{D}{$\displaystyle\int_0^{\sqrt2+1}\dfrac{t}{(1+t^2)\sqrt{1-t^2}}\,dt$}
\end{qoptions}

\qn{12} The value of $\displaystyle\int_0^{e-1}\dfrac{e^{\frac{x^2+2x-1}{2}}}{(x+1)}\,dx+\int_1^e x.\ell nx\,e^{\frac{x^2-2}{2}}\,dx$ is
\begin{qoptions}
\opt{A}{$\sqrt e^{\,(e^2+1)}$}
\opt{B}{$\sqrt e^{\,(e^2-1)}$}
\opt{C}{0}
\opt{D}{$\left(\sqrt e\right)^{e^2-2}$}
\end{qoptions}

\qn{13} A water tank has the shape of an inverted circular cone with its axis vertical and vertex lowermost. Its semi-vertical angle is $\tan^{-1}(0.5)$. Water is poured into it at a constant rate of 5 cubic meter per hour. The rate at which the level of water is rising at the instant when the depth of water in tank is 4m is $\left(\text{Take }\pi=\dfrac{22}7\right)$
\begin{qoptions}
\opt{A}{1m/h}
\opt{B}{$\dfrac{30}{17}$ m/h}
\opt{C}{$\dfrac{70}{88}$ m/h}
\opt{D}{$\dfrac{35}{88}$ m/h}
\end{qoptions}

\qn{14} Let $f(x)=x^2-x\sin x-\cos x$, then which of following is \underline{false} ?
\begin{qoptions}
\opt{A}{$f(x)=k$ has no solution for $k<-1$}
\opt{B}{$f(x)$ is decreasing for $x<0$ and increasing for $x>0$}
\opt{C}{$\displaystyle\lim_{x\to\infty}f(x)$ is not finite}
\opt{D}{The roots of $f(x)=0$ lie on same side of origin}
\end{qoptions}

\qn{15} \textbf{STATEMENT 1 :} A number is chosen at random from the numbers 1, 2, 3, \dots (6n + 3). Let events A and B be defined as follows\\
A : number is divisible by 2\\
B : number is divisible by 3\\
Then, A and B are independent.\\
\textbf{STATEMENT 2 :} If events A and B are independent, then $P(A\cap B)=P(A).P(B)$.
\begin{qoptions}
\opt{A}{Both statements are true}
\opt{B}{Both statements are false}
\opt{C}{Statement - 1 is true and statement - 2 is false}
\opt{D}{Statement - 1 is false and statement - 2 is true}
\end{qoptions}

\qn{16} Value of $\sin^{-1}\left(\dfrac{x^2}4+\dfrac{y^2}9\right)$\\$+\cos^{-1}\left(\dfrac{x}{2\sqrt2}+\dfrac{y}{3\sqrt2}-2\right)$ is
\begin{qoptions}
\opt{A}{$\dfrac\pi2$}
\opt{B}{$\pi$}
\opt{C}{$\dfrac{3\pi}2$}
\opt{D}{$2\pi$}
\end{qoptions}

\qn{17} If the integral $\displaystyle\int\dfrac{5\tan x\,dx}{\tan x-2}=x+a\,\ell n|\sin x-2\cos x|+C$ then `$a$' is equal to ($C$ is constant of integration)
\begin{qoptions}
\opt{A}{1}
\opt{B}{2}
\opt{C}{$-1$}
\opt{D}{$-2$}
\end{qoptions}

\qn{18} Let $(x^2+y^2)dy=xy\,dx$ and $y(1)=1$. If $y(\alpha)=e$, then $\left(\dfrac\alpha e\right)^2$ is
\begin{qoptions}
\opt{A}{$\sqrt3$}
\opt{B}{3}
\opt{C}{$3\sqrt3$}
\opt{D}{9}
\end{qoptions}

\qn{19} Let the lines $\dfrac{x-1}{\lambda}=\dfrac{y-2}{1}=\dfrac{z-3}{2}$ and $\dfrac{x+26}{-2}=\dfrac{y+18}{3}=\dfrac{z+28}{\lambda}$ be coplanar and P be the plane containing these two lines. Then which of the following points does not lie on P ?
\begin{qoptions}
\opt{A}{$(0,-2,-2)$}
\opt{B}{$(-5,0,-1)$}
\opt{C}{$(3,-1,0)$}
\opt{D}{$(0,4,5)$}
\end{qoptions}

\qn{20} If $x=\tan\dfrac\pi8$, then the value of $x^4+3x^3-3x^2-9x+8$ is
\begin{qoptions}
\opt{A}{4}
\opt{B}{5}
\opt{C}{6}
\opt{D}{7}
\end{qoptions}

\sect{Section-II}

\qn{1} Find number of possible values of $\beta$, if $\sqrt{2\cot\alpha+\dfrac1{\sin^2\alpha}}=\dfrac{1+\cot\alpha}{\sqrt{\tan\beta}}$ for $\alpha\in\left(\dfrac{3\pi}4,\pi\right)$

\qn{2} Find the integer nearest to the maximum possible value of `$m$', such that the area of triangle formed by the lines $y=m(x+1)$, $y=x$ and $y=-x$ is 4 square units.

\qn{3} If the sum of rational terms in the expansion of $\left(\sqrt[3]2+\sqrt[6]{\dfrac13}\right)^6$ is `$a$' and coefficient of term independent of $x$ in $(x^{-1}+x)^6$ is `$b$', then the eccentricity of the ellipse whose major and minor axes are $b$ and $a$ respectively is $e$. The value of $[10e]$, where $[\cdot]$ denotes G.I.F. is

\qn{4} If the system of equations $2x-3y+5z=12$, $3x+y+\lambda z=\mu$, $x-7y+8z=17$ has infinitely many real solutions, then $\lambda+\mu=$ .......

\qn{5} Find the total number of ways in which a 3-digit number $xyz$ can be formed from the digits 1, 2, 3, 4, 5, 6, 7, 8 and 9 such that $x,y,z$ are in A.P and $x\neq y\neq z$.

\qn{6} If the eccentricity of the ellipse $\dfrac{x^2}{a^2+1}+\dfrac{y^2}{a^2+2}=1$ be $\dfrac1{\sqrt6}$, then the integer nearest to the latus rectum of the ellipse is .........

\qn{7} Let $S=\{4,6,9\}$ and $T=\{9,10,11,\dots,1000\}$. If $A=\{a_1+a_2+\dots+a_k:k\in N,\ a_1,a_2,a_3,\dots,a_k\in S\}$, then the sum of all the elements in the set $T-A$ is equal to \underline{\hspace{3em}}.

\qn{8} If the domain of the function $f(x)=\log_e(4x^2+11x+6)+\sin^{-1}(4x+3)+\cos^{-1}\left(\dfrac{10x+6}3\right)$ is $(\alpha,\beta]$, then $36|\alpha+\beta|$ is equal to

\qn{9} Let $m_1$ and $m_2$ be the slopes of the tangents drawn from the point P(4, 1) to the hyperbola H : $\dfrac{y^2}{25}-\dfrac{x^2}{16}=1$. If Q is the point form which the tangents drawn to H have slopes $|m_1|$ and $|m_2|$ and they make positive intercepts $\alpha$ and $\beta$ on the $x$-axis, then $\dfrac{(PQ)^2}{\alpha\beta}$ is equal to

\qn{10} Nine nonnegative number have average 10. What is the greatest possible value for their median ?

\end{multicols}
\end{document}
